% collision_asymmetry: why momentum survives a collision and kinetic energy need not (Ch 12, fig:collasym)
% Left: the third-law pair during contact.  Body 1 pushes body 2 with F_12 (drawn on
% body 2, to the right); body 2 pushes back on body 1 with F_21 = -F_12 (drawn on
% body 1, to the left).  Equal lengths: the pair cancels in dP/dt.
% Right: the same pair squashed at the contact; the strokes mark the kinetic energy
% that leaves as heat, sound and permanent deformation.
\begin{tikzpicture}[x=1cm,y=1cm]
  \def\r{0.5}
  % ---------------------------------------------------------------- momentum
  \begin{scope}
    \node[note,text=cTeal] at (2.6,1.25) {momentum is conserved};
    \draw[body] (2.1,0) circle[radius=\r];
    \draw[body,fill=cMoon] (3.1,0) circle[radius=\r];
    \node at (2.1,0.22) {1};
    \node at (3.1,0.22) {2};
    \draw[force] (2.1,-0.12) -- ++(-1.25,0) node[anchor=east] {$\vF_{21}$};
    \draw[force] (3.1,-0.12) -- ++(1.25,0) node[anchor=west] {$\vF_{12}$};
    \node[font=\footnotesize,align=center] at (2.6,-1.15)
      {equal and opposite at every instant,\\ so the total momentum cannot change};
  \end{scope}
  % ---------------------------------------------------------- kinetic energy
  \begin{scope}[xshift=6.7cm]
    \node[note,text=cAcc] at (2.6,1.25) {kinetic energy may be lost};
    % two balls flattened where they press: centres 0.86 apart, flat at the midpoint
    \pgfmathsetmacro\ang{acos(0.43/\r)}
    \draw[body] (2.6,{\r*sin(\ang)}) arc[start angle=\ang,end angle=360-\ang,radius=\r] -- cycle;
    \draw[body,fill=cMoon] (2.6,{-\r*sin(\ang)}) arc[start angle=180+\ang,end angle=540-\ang,radius=\r] -- cycle;
    \node at (2.17,0.22) {1};
    \node at (3.03,0.22) {2};
    % the bang: strokes radiating from the contact
    \foreach \t in {65,90,115,245,270,295}
      \draw[cAcc,thin line,line cap=round] (2.6,0) ++(\t:0.6) -- ++(\t:0.3);
    \node[cAcc,font=\footnotesize] at (2.6,-1.13) {heat, sound, deformation};
    \node[font=\footnotesize] at (2.6,-1.5) {only an elastic collision keeps it all};
  \end{scope}
\end{tikzpicture}
